\begin{theorem}[MetaCIC candidate-mediated SN obligation surface]
\label{thm:metacic-confluence-candidate-mediated-sn-obligation-surface}
Every confluence-facing or decidability-facing consumer that uses candidate-mediated strong-normalization evidence from an accepted $\MetaCICCriticalPathUp$ packet must first display the strong-normalization row $S$, the normal-form endpoint $N$, the handoff row $U$, the discharge socket $D$, the componentwise transport row $H$, the displayed continuation row $C$, the provenance row $P$, the local name $L$, and the four L10 faces
$$
\DyadicRatCoreUp,\quad \StreamNameUp,\quad \RegSeqRatUp,\quad \RealUp
$$
before it reads any confluence or decidability budget. The consumer may inspect only this candidate-mediated surface and the existing handoff budget; it receives no global confluence theorem, no beta-conversion decision procedure, no filled discharge socket, no quotient term equality, no host real equality, no selected limit, and no $\RealUp$ completeness claim.
\end{theorem}
\begin{proof}
Begin with the candidate-mediated SN frontier, \autoref{thm:metacic-critical-path-candidate-mediated-sn-frontier}. It confines the strong-normalization-facing read to the displayed frontier rows and requires the four L10 faces before any confluence-facing or decidability-facing consumer reads the handoff row.

Next apply the confluence-decidability budget, \autoref{thm:metaciccriticalpath-confluence-decidability-budget}. That budget keeps every confluence or decidability read downstream of the open-source ledger and the finite packet route, so the consumer cannot read $U$ as an independent global theorem.

Finally use the consumer boundary, \autoref{thm:metacic-critical-path-confluence-decidability-consumer-boundary}. It allows closed-normal-form membership to pass through the handoff route while excluding closed subject reduction, full consistency, arbitrary Church--Rosser, global beta-conversion decidability, discharged sockets, quotient equality, host equality, selected limits, and $\RealUp$ completeness. Therefore the listed rows form the whole obligation surface available before any confluence or decidability budget is consumed.
\end{proof}

\begin{theorem}[MetaCIC candidate-discharge schedule exhaustion]
\label{thm:metacic-critical-path-candidate-discharge-schedule-exhaustion}
Any accepted $\MetaCICCriticalPathUp$ consumer that attempts to discharge a candidate-mediated strong-normalization or confluence socket must expose the candidate row $K$, the strong-normalization row $S$, the normal-form endpoint $N$, the confluence-decidability handoff row $U$, the discharge socket $D$, the four L10 faces
$$
\DyadicRatCoreUp,\quad \StreamNameUp,\quad \RegSeqRatUp,\quad \RealUp,
$$
componentwise transport $H$, displayed replay $C$, provenance $P$, and the local-name row $L$ before it reads a confluence or decidability budget. The discharge schedule exports no closed subject-reduction theorem, arbitrary Church--Rosser theorem, global beta-conversion decision procedure, filled discharge socket, quotient term equality, selected limit, host real equality, or $\RealUp$ completeness claim.
\end{theorem}
\begin{proof}
Begin with the candidate-mediated obligation surface, \autoref{thm:metacic-confluence-candidate-mediated-sn-obligation-surface}. It already forces $S,N,U,D,H,C,P,L$ and the four L10 faces to be displayed before any confluence-facing or decidability-facing budget can be read.

Now add the candidate row $K$ as the schedule witness being consumed. The candidate-mediated frontier and source-ledger route keep this row attached to the same accepted $\MetaCICCriticalPathUp$ packet rather than letting it appear as a detached normalization theorem.

Finally apply the confluence-decidability budget \autoref{thm:metaciccriticalpath-confluence-decidability-budget} and the consumer boundary \autoref{thm:metacic-critical-path-confluence-decidability-consumer-boundary}. Those results allow the downstream budget read only after the displayed packet route is present and exclude closed subject reduction, full consistency, arbitrary Church--Rosser, global beta-conversion decidability, discharged sockets, quotient equality, host equality, selected limits, and $\RealUp$ completeness. Thus the listed rows exhaust the candidate-discharge schedule before any confluence or decidability budget is consumed.
\end{proof}
\leanchecked{BEDC.Derived.MetaCICCriticalPathUp.MetaCICCriticalPathCandidateDischargeScheduleExhaustion}

\begin{theorem}[MetaCIC candidate-mediated SN discharge]
\label{thm:metacic-critical-path-candidate-mediated-sn-discharge}
Any accepted $\MetaCICCriticalPathUp$ consumer that routes candidate-mediated strong-normalization evidence into a confluence-facing or decidability-facing budget must keep the discharge socket $D$ explicitly displayed and unfilled throughout the read. The allowed route is the finite surface consisting of the candidate row, the strong-normalization row $S$, the normal-form endpoint $N$, the handoff row $U$, componentwise transport $H$, displayed replay $C$, provenance $P$, local naming $L$, and the four L10 faces
$$
\DyadicRatCoreUp,\quad \StreamNameUp,\quad \RegSeqRatUp,\quad \RealUp .
$$
The route may feed the displayed confluence/decidability budget only as a conditional socket surface; it exports no closed subject-reduction row, no filled discharge socket, no global confluence theorem, no beta-conversion decision procedure, no quotient term equality, no selected limit, and no ambient $\RealUp$ completeness claim.
\end{theorem}
\leanchecked{BEDC.Derived.MetaCICCriticalPathUp.MetaCICCriticalPathCandidateMediatedSNDischarge}
\begin{proof}
First apply the candidate-mediated obligation surface, \autoref{thm:metacic-confluence-candidate-mediated-sn-obligation-surface}. It makes $S,N,U,D,H,C,P,L$ and the four L10 faces visible before any downstream budget can be inspected, so the socket row $D$ is part of the source surface rather than an inferred conclusion.

Next use the candidate-discharge schedule exhaustion theorem, \autoref{thm:metacic-critical-path-candidate-discharge-schedule-exhaustion}. This adds the candidate row and shows that every confluence-facing or decidability-facing read must consume the displayed schedule before touching the budget rows.

Finally apply the frontier discharge row, \autoref{thm:metacic-critical-path-candidate-mediated-sn-discharge-row}. That theorem keeps $D$ separate from the candidate-mediated SN evidence while the finite frontier is replayed. Since no row in the displayed route is a closed subject-reduction coordinate, the consumer receives only the conditional discharge surface and none of the excluded global MetaCIC, quotient, selected-limit, or completeness claims.
\end{proof}

\begin{theorem}[MetaCIC candidate-mediated SN conditional budget]
\label{thm:metacic-critical-path-candidate-mediated-sn-conditional-budget}
Any accepted $\MetaCICCriticalPathUp$ consumer may pass candidate-mediated strong-normalization evidence to a confluence-facing or decidability-facing budget only when the candidate row $K$, strong-normalization row $S$, normal-form endpoint $N$, handoff row $U$, unfilled discharge socket $D$, componentwise transport $H$, displayed replay $C$, provenance $P$, local name $L$, and the four L10 faces
$$
\DyadicRatCoreUp,\quad \StreamNameUp,\quad \RegSeqRatUp,\quad \RealUp
$$
are all displayed. The conclusion remains conditional on the socket row $D$ and exports no closed subject-reduction theorem, no global Church--Rosser theorem, no beta-conversion decision procedure, no quotient term equality, no selected limit, and no $\RealUp$ completeness claim.
\end{theorem}
\begin{proof}
First use the obligation surface of \autoref{thm:metacic-confluence-candidate-mediated-sn-obligation-surface}. It exposes the strong-normalization row, normal-form endpoint, handoff row, discharge socket, transport, replay, provenance, local name, and the four L10 faces before any downstream budget can be read.

Next apply the schedule exhaustion theorem, \autoref{thm:metacic-critical-path-candidate-discharge-schedule-exhaustion}. It adds the candidate row $K$ and shows that every candidate-mediated route must display the same finite schedule before touching confluence-facing or decidability-facing data.

Finally apply the discharge theorem, \autoref{thm:metacic-critical-path-candidate-mediated-sn-discharge}. It keeps the socket unfilled while the route is replayed, so the consumer receives only a conditional budget and none of the excluded closed MetaCIC, quotient, selected-limit, or Real-completeness claims.
\end{proof}
\leanchecked{BEDC.Derived.MetaCICCriticalPathUp.MetaCICCriticalPathCandidateMediatedSNConditionalBudget}

\begin{theorem}[MetaCIC four-face bounded checker gate]
\label{thm:metacic-critical-path-four-face-bounded-checker-gate}
Any bounded checker consumer that uses candidate-mediated strong-normalization evidence from an accepted $\MetaCICCriticalPathUp$ packet must display the candidate row $K$, strong-normalization row $S$, normal endpoint row $N$, handoff row $U$, open discharge socket $D$, componentwise transport $H$, displayed replay $C$, provenance $P$, local naming $L$, and the four L10 faces
$$
\DyadicRatCoreUp,\quad \StreamNameUp,\quad \RegSeqRatUp,\quad \RealUp
$$
before it reads the checker budget. The gate exports no closed subject-reduction theorem, global confluence theorem, beta-conversion decision procedure, quotient term equality, selected limit, host equality, or $\RealUp$ completeness claim.
\end{theorem}
\leanchecked{BEDC.Derived.MetaCICCriticalPathUp.MetaCICCriticalPathFourFaceBoundedCheckerGate}
\begin{proof}
Begin with \autoref{thm:metacic-confluence-candidate-mediated-sn-obligation-surface}. It places the strong-normalization row, normal endpoint, handoff row, discharge socket, transport, replay, provenance, local name, and four L10 faces before any confluence-facing or decidability-facing budget.

Apply \autoref{thm:metacic-criticalpath-candidate-mediated-sn-premise-split}. The candidate row $K$, the $S,N$ endpoint row, the open socket $D$, and the L10 face row remain separate premises rather than a single closed normalization theorem.

Use \autoref{thm:metacic-critical-path-candidate-mediated-sn-conditional-budget}. The checker budget may be read only after the displayed schedule is present, and the conclusion remains conditional on the socket row.

Finally apply \autoref{thm:metacic-critical-path-bounded-normalization-candidate-ledger}. The bounded-normalization candidate ledger supplies checker provenance for the same candidate-mediated route, while the accepted packet still has no coordinate for subject reduction, global confluence, beta-conversion decidability, quotient term equality, selected limits, host equality, or Real completeness.
\end{proof}

\begin{theorem}[MetaCIC bounded checker frontier separation]
\label{thm:metacic-critical-path-bounded-checker-frontier-separation}
Any bounded checker read that consumes candidate-mediated strong-normalization evidence from an accepted $\MetaCICCriticalPathUp$ packet must keep the candidate row $K$, strong-normalization endpoint rows $S,N$, open discharge socket $D$, and the four L10 faces
$$
\DyadicRatCoreUp,\quad \StreamNameUp,\quad \RegSeqRatUp,\quad \RealUp
$$
as separate displayed premises before the checker budget is read. The checker may pass those premises through the displayed handoff, transport, replay, provenance, and local naming rows, but it may not collapse them into subject reduction, global confluence, beta-conversion decidability, quotient equality, selected limit, or $\RealUp$ completeness.
\end{theorem}
\begin{proof}
Start with the obligation surface of \autoref{thm:metacic-confluence-candidate-mediated-sn-obligation-surface}. It displays the SN endpoint, discharge socket, handoff, transport, replay, provenance, local name, and four L10 faces before any checker-facing budget can be consumed.

Apply \autoref{thm:metacic-critical-path-candidate-discharge-schedule-exhaustion}. The candidate row $K$ is part of the finite schedule, so it must remain visible before the bounded checker reads its budget.

Use \autoref{thm:metacic-critical-path-candidate-mediated-sn-conditional-budget} and \autoref{thm:metacic-critical-path-four-face-bounded-checker-gate}. These keep the route conditional on $D$ and require the four L10 faces to be displayed as the checker boundary.

Finally apply \autoref{thm:metacic-criticalpath-candidate-mediated-sn-premise-split}. The candidate row, endpoint row, socket row, and L10 face row remain separate premises, so none can be repacked as closed subject reduction, global confluence, beta-conversion decidability, quotient equality, a selected limit, or Real completeness.
\end{proof}

\leanchecked{BEDC.Derived.MetaCICCriticalPathUp.MetaCICCriticalPathBoundedCheckerFrontierSeparation}

\begin{theorem}[MetaCIC bounded checker budget factorization]
\label{thm:metacic-critical-path-bounded-checker-budget-factorization}
Any bounded checker budget read that consumes candidate-mediated strong-normalization evidence from an accepted $\MetaCICCriticalPathUp$ packet factors through four separate displayed premise rows:
\begin{enumerate}
\item the candidate row $K$;
\item the strong-normalization endpoint row $S,N$;
\item the open discharge socket row $D$;
\item the four L10 faces
$$
\DyadicRatCoreUp,\quad \StreamNameUp,\quad \RegSeqRatUp,\quad \RealUp .
$$
\end{enumerate}
Only after these rows are displayed may the checker read the handoff, transport, replay, provenance, local naming, and bounded budget rows. The factorization exports no closed subject-reduction theorem, global confluence theorem, beta-conversion decision procedure, quotient term equality, selected limit, host equality, or $\RealUp$ completeness claim.
\end{theorem}
\begin{proof}
Begin with \autoref{thm:metacic-critical-path-four-face-bounded-checker-gate}. It requires the candidate row, SN endpoint rows, open socket, and four L10 faces to be displayed before the bounded checker reads its budget.

Next apply \autoref{thm:metacic-critical-path-candidate-mediated-sn-conditional-budget}. The checker-facing budget remains conditional on the socket row $D$, so the budget cannot be read as a closed normalization theorem.

Use \autoref{thm:metacic-criticalpath-candidate-mediated-sn-premise-split}. It separates $K$, the endpoint row $S,N$, the socket row $D$, and the L10 face row from the handoff and structural rows.

Finally read the frontier separation theorem, \autoref{thm:metacic-critical-path-bounded-checker-frontier-separation}. The same separation persists at the bounded-checker frontier, so the checker receives a factored finite budget rather than subject reduction, global confluence, beta-decision, quotient equality, selected-limit, host-equality, or Real-completeness data.
\end{proof}

\leanchecked{BEDC.Derived.MetaCICCriticalPathUp.MetaCICCriticalPathBoundedCheckerBudgetFactorization}

\begin{theorem}[MetaCIC bounded checker replay separation]
\label{thm:metacic-critical-path-bounded-checker-replay-separation}
Every replay of a bounded-checker budget read from an accepted $\MetaCICCriticalPathUp$ packet must preserve four separate premise rows throughout the replay:
\begin{enumerate}
\item the candidate row $K$;
\item the strong-normalization endpoint row $S,N$;
\item the open discharge socket row $D$;
\item the L10 face row over $\DyadicRatCoreUp$, $\StreamNameUp$, $\RegSeqRatUp$, and $\RealUp$.
\end{enumerate}
The replay may then read the handoff, transport, displayed continuation, provenance, local naming, and bounded budget rows, but it cannot compress the four premise rows into closed subject reduction, global confluence, beta-conversion decidability, quotient equality, selected limits, host equality, or $\RealUp$ completeness.
\end{theorem}
\leanchecked{BEDC.Derived.MetaCICCriticalPathUp.MetaCICCriticalPathBoundedCheckerReplaySeparation}
\begin{proof}
Start with the candidate-mediated obligation surface of \autoref{thm:metacic-confluence-candidate-mediated-sn-obligation-surface}. It exposes the endpoint, socket, handoff, transport, replay, provenance, local name, and four L10 faces before any checker-facing budget is read.

Use schedule exhaustion from \autoref{thm:metacic-critical-path-candidate-discharge-schedule-exhaustion}. The candidate row remains attached to the displayed finite route, so replay cannot detach $K$ from the same accepted packet.

Apply the conditional budget theorem, \autoref{thm:metacic-critical-path-candidate-mediated-sn-conditional-budget}, together with the four-face bounded checker gate, \autoref{thm:metacic-critical-path-four-face-bounded-checker-gate}. These rows keep the socket open and require the dyadic, stream, regular-rational, and real faces to appear as checker premises.

Finally apply \autoref{thm:metacic-critical-path-bounded-checker-budget-factorization}. Its factorization is stable under replay because the replay rows are read after the four premise rows, so the bounded checker receives only the displayed conditional budget and none of the excluded closed MetaCIC, quotient, selected-limit, host-equality, or completeness claims.
\end{proof}

\begin{theorem}[MetaCIC candidate-mediated SN frontier coverage]
\label{thm:metaciccriticalpath-candidate-mediated-sn-frontier-coverage}
Any accepted $\MetaCICCriticalPathUp$ consumer crossing from candidate-mediated strong-normalization evidence into confluence-facing or decidability-facing use must display the whole frontier before the handoff is read:
\begin{enumerate}
\item the candidate row $K$;
\item the strong-normalization endpoint and normal-form endpoint rows $S,N$;
\item the open discharge socket row $D$;
\item the handoff row $U$;
\item the bounded checker budget row;
\item the four L10 faces
$$
\DyadicRatCoreUp,\quad \StreamNameUp,\quad \RegSeqRatUp,\quad \RealUp .
$$
\end{enumerate}
The frontier may then pass through componentwise $\hsame$ transport, displayed $\Cont$ replay, $\Pkg$ provenance, and the local naming row. It does not export subject reduction, global confluence, beta-conversion decidability, quotient term equality, selected limits, host equality, or $\RealUp$ completeness.
\end{theorem}
\begin{proof}
First read the obligation surface of \autoref{thm:metacic-confluence-candidate-mediated-sn-obligation-surface}. It exposes the endpoint rows, open socket, handoff, transport, replay, provenance, local name, and four L10 faces before any confluence-facing or decidability-facing consumer reads the route.

Next apply \autoref{thm:metacic-critical-path-candidate-mediated-sn-conditional-budget}. The bounded checker budget remains conditional on the displayed candidate-mediated surface, so it cannot be detached from $K$, $S,N$, $D$, or $U$.

Use \autoref{thm:metacic-critical-path-four-face-bounded-checker-gate}. The dyadic, stream, regular-rational, and real faces must be displayed as four separate L10 premises, not compressed into a single ambient Real-completeness assumption.

Finally combine \autoref{thm:metacic-critical-path-bounded-checker-budget-factorization} with the frontier separation theorem. The factorization keeps the bounded budget behind the same frontier rows, while separation prevents the consumer from repacking the finite route as subject reduction, global confluence, beta-conversion decidability, quotient term equality, selected limits, host equality, or Real completeness.
\end{proof}
\leanchecked{BEDC.Derived.MetaCICCriticalPathUp.MetaCICCriticalPathCandidateMediatedSNFrontierCoverage}

\begin{theorem}[MetaCIC candidate-mediated SN premise split]
\label{thm:metacic-criticalpath-candidate-mediated-sn-premise-split}
Any consumer passing candidate-mediated strong-normalization evidence from an accepted $\MetaCICCriticalPathUp$ packet into a confluence-facing or decidability-facing budget must display four separate premise rows before the budget is read:
\begin{enumerate}
\item the candidate row $K$;
\item the strong-normalization and normal-form endpoint row consisting of $S$ and $N$;
\item the unfilled discharge socket row $D$;
\item the L10 face row over $\DyadicRatCoreUp$, $\StreamNameUp$, $\RegSeqRatUp$, and $\RealUp$.
\end{enumerate}
The four rows remain separate from the handoff row $U$, componentwise transport $H$, displayed replay $C$, provenance $P$, and local name $L$. The consumer receives no closed subject-reduction theorem, arbitrary Church--Rosser theorem, global beta-conversion decision, quotient term equality, selected limit, host real equality, or Real completeness.
\end{theorem}
\begin{proof}
Start with \autoref{thm:metacic-confluence-candidate-mediated-sn-obligation-surface}. It places the strong-normalization row, normal-form endpoint, handoff row, discharge socket, transport, replay, provenance, local name, and L10 faces before any confluence-facing or decidability-facing budget.

Next apply \autoref{thm:metacic-critical-path-candidate-discharge-schedule-exhaustion}. This adds the candidate row $K$ as a displayed schedule premise and keeps it attached to the same accepted packet rather than to a detached normalization theorem.

Use \autoref{thm:metacic-critical-path-candidate-mediated-sn-discharge}. The discharge socket $D$ remains visible and unfilled while the candidate-mediated evidence is routed, so the socket is not absorbed into the strong-normalization endpoint.

Finally apply \autoref{thm:metacic-critical-path-candidate-mediated-sn-conditional-budget}. The budget is conditional on all displayed rows, and the L10 face row remains the finite dependency surface over $\DyadicRatCoreUp$, $\StreamNameUp$, $\RegSeqRatUp$, and $\RealUp$. Since the accepted route has no closed subject-reduction coordinate, arbitrary Church--Rosser theorem, global beta-conversion decision, quotient equality, selected limit, host equality, or Real-completeness row, none of those principles is exported.
\end{proof}

\leanchecked{BEDC.Derived.MetaCICCriticalPathUp.MetaCICCriticalPathCandidateMediatedSNObligationSurface}
\leanchecked{BEDC.Derived.MetaCICCriticalPathUp.MetaCICCriticalPathCandidateMediatedSNPremiseSplit}

\begin{theorem}[MetaCIC critical-path bounded conversion discharge]
\label{thm:metacic-critical-path-bounded-conversion-discharge}
Every bounded-conversion consumer of an accepted $\MetaCICCriticalPathUp$ packet must consume the displayed $\CandidateMediatedSNUp$ row, the $\MetaCICNormalizationFrontierUp$ boundary row, the $\MetaCICCriticalPairLedgerUp$ pair ledger, and the L10 four-object exit certificate rows over $\DyadicRatCoreUp$, $\StreamNameUp$, $\RegSeqRatUp$, and $\RealUp$ before reading any conversion budget. The discharge route also carries the endpoint row $S,N$, the open socket row $D$, componentwise $\hsame$ transport, displayed $\Cont$ replay, $\Pkg$ provenance, and the local naming row.

The discharge is finite and bounded by the displayed packet. It does not assert global strong normalization, arbitrary Church--Rosser, unrestricted beta-conversion decision, quotient term equality, selected limits, host real equality, or ambient $\RealUp$ completeness.
\end{theorem}
\begin{proof}
Start with \autoref{thm:metacic-criticalpath-candidate-mediated-sn-premise-split}. It requires the candidate row, endpoint row, socket row, and L10 face row to be displayed separately before the budget is read.

Apply \autoref{thm:metacic-critical-path-bounded-checker-frontier-separation}. The normalization frontier remains a boundary row rather than a closed normalization theorem, so the bounded checker receives only the exposed finite interface.

Use \autoref{thm:metacic-critical-path-candidate-discharge-schedule-exhaustion} together with the critical-pair ledger row. The candidate-mediated schedule and pair ledger are consumed as finite discharge premises, while $H,C,P,L$ transport, replay, source, and name the same route.

Since the socket remains visible and the L10 faces remain explicit, the route cannot be repacked as global strong normalization, unrestricted conversion decidability, quotient equality, selected-limit data, host equality, or Real completeness.
\end{proof}
\leanchecked{BEDC.Derived.MetaCICCriticalPathUp.MetaCICCriticalPathBoundedConversionDischarge}

\begin{theorem}[MetaCIC bounded conversion socket budget]
\label{thm:metacic-critical-path-bounded-conversion-socket-budget}
Every bounded-conversion consumer of an accepted $\MetaCICCriticalPathUp$ packet must display the candidate row, the strong-normalization endpoint rows $S,N$, the open discharge socket $D$, the bounded checker budget row, and the four L10 faces
$$
\DyadicRatCoreUp,\quad \StreamNameUp,\quad \RegSeqRatUp,\quad \RealUp
$$
before it reads any conversion data. The socket remains open throughout the read, and the packet exports no subject reduction, global Church--Rosser theorem, beta-decision procedure, quotient equality, selected limit, host equality, or $\RealUp$ completeness.
\end{theorem}
\leanchecked{BEDC.Derived.MetaCICCriticalPathUp.MetaCICCriticalPathBoundedConversionSocketBudget}
\begin{proof}
Start with the candidate-mediated obligation surface, \autoref{thm:metacic-confluence-candidate-mediated-sn-obligation-surface}. It places the endpoint rows, open socket, structural rows, and the four L10 faces before any conversion-facing budget can be inspected.

Apply the candidate-discharge schedule exhaustion, \autoref{thm:metacic-critical-path-candidate-discharge-schedule-exhaustion}. The candidate row must remain attached to the same accepted packet, so the bounded-conversion consumer cannot read it as a detached normalization theorem.

Use the bounded checker budget factorization, \autoref{thm:metacic-critical-path-bounded-checker-budget-factorization}. It separates the candidate row, endpoint rows, socket row, and L10 face row before the checker budget is read.

Finally apply the bounded conversion discharge theorem, \autoref{thm:metacic-critical-path-bounded-conversion-discharge}. The conversion read is finite and conditional on the displayed socket, so none of the excluded subject-reduction, Church--Rosser, beta-decision, quotient-equality, selected-limit, host-equality, or Real-completeness data is exported.
\end{proof}

\begin{theorem}[MetaCIC critical-path candidate-SN readiness]
\label{thm:metacic-critical-path-candidate-sn-readiness}
An accepted $\MetaCICCriticalPathUp$ packet is ready for candidate-mediated strong-normalization handoff precisely at the finite surface consisting of the $\CandidateMediatedSNUp$ row, the $\MetaCICNormalizationFrontierUp$ frontier row, the $\MetaCICCriticalPairLedgerUp$ ledger row, the endpoint row $S,N$, the open discharge socket $D$, the L10 four-object exit certificate rows, componentwise $\hsame$ transport, displayed $\Cont$ replay, $\Pkg$ provenance, and the local naming row.

Readiness means only that those rows are all present and factor the bounded-conversion consumer route. It does not close subject reduction, confluence, global strong normalization, beta-conversion decidability, quotient term equality, selected limits, host equality, or Real completeness.
\end{theorem}
\begin{proof}
Read the candidate-mediated obligation surface and the premise split. Together they display the candidate row, endpoint row, socket row, handoff row, structural rows, and L10 faces before any confluence-facing or decidability-facing budget is consumed.

The normalization frontier and critical-pair ledger are then read as finite boundary rows. They record where bounded conversion may discharge local obligations, not a global theorem that all conversions terminate or join.

Finally apply \autoref{thm:metacic-critical-path-bounded-conversion-discharge}. The same displayed rows are exactly the finite discharge interface for bounded conversion, so candidate-SN readiness is a row-presence and factorization claim only.
\end{proof}
\leanchecked{BEDC.Derived.MetaCICCriticalPathUp.MetaCICCriticalPathCandidateSNReadiness}

\begin{theorem}[MetaCIC residual candidate-SN handoff]
\label{thm:metacic-critical-path-residual-candidate-sn-handoff}
An accepted $\MetaCICCriticalPathUp$ packet may pass candidate-mediated strong-normalization evidence to residual-substitution or confluence consumers only after it displays the $\CandidateMediatedSNUp$ row, the $\MetaCICNormalizationFrontierUp$ row, the $\MetaCICResidualSubstitutionBoundaryUp$ row, the $\MetaCICCriticalPairLedgerUp$ row, the endpoint rows $S,N$, the open discharge socket $D$, the four L10 faces $\DyadicRatCoreUp$, $\StreamNameUp$, $\RegSeqRatUp$, and $\RealUp$, and the structural transport, replay, provenance, and local-name rows.

The handoff is a finite consumer boundary. It permits residual-substitution and confluence consumers to read the candidate-mediated evidence only through the displayed socket and bounded-conversion discharge rows, and it does not export subject reduction, unrestricted confluence, global strong normalization, beta-conversion decidability, quotient term equality, selected limits, host equality, or ambient $\RealUp$ completeness.
\end{theorem}
\leanchecked{BEDC.Derived.MetaCICCriticalPathUp.MetaCICCriticalPathResidualCandidateSNHandoff}
\begin{proof}
Start with the candidate-mediated obligation surface. It places the candidate row, endpoint rows, open socket, structural rows, and the four L10 faces before any normalization or conversion consumer can inspect the evidence.

Next apply the candidate-SN readiness theorem. Readiness supplies the finite row-presence and factorization condition, so the evidence remains attached to the accepted $\MetaCICCriticalPathUp$ packet rather than becoming a detached strong-normalization theorem.

Use the premise split and bounded-conversion discharge rows. The premise split separates the candidate row from residual-substitution, frontier, pair-ledger, endpoint, socket, and L10 obligations, while the bounded-conversion discharge keeps the consumer route conditional on those displayed rows.

Finally repack through the structural coordinates. Componentwise $\hsame$ transport, displayed $\Cont$ replay, $\Pkg$ provenance, and the local naming row preserve the same packet boundary; hence residual-substitution and confluence consumers receive only the displayed candidate-mediated handoff surface.
\end{proof}

\begin{theorem}[MetaCIC critical-path bounded conversion consumer route]
\label{thm:metacic-critical-path-bounded-conversion-consumer-route}
Any confluence-facing or decidability-facing consumer that uses the accepted candidate-mediated strong-normalization packet of \autoref{thm:metacic-critical-path-candidate-sn-readiness} must first read the bounded normalization candidates, the candidate discharge schedules, the finite checker budgets, and the local $\NameCert_{\MetaCICCriticalPathUp}$ rows. Only after these rows are displayed may the consumer export a bounded conversion claim through the $\MetaCICNormalizationFrontierUp$ boundary, the $\MetaCICCriticalPairLedgerUp$ pair ledger, the endpoint rows $S,N$, the open discharge socket $D$, and the L10 four-object exit certificate over $\DyadicRatCoreUp$, $\StreamNameUp$, $\RegSeqRatUp$, and $\RealUp$.

The route is a consumer discipline, not a global conversion theorem. It does not assert unrestricted confluence, global strong normalization, unbounded beta-conversion decidability, quotient term equality, selected limits, host equality, or ambient $\RealUp$ completeness.
\end{theorem}
\begin{proof}
Begin with the candidate-mediated SN readiness surface. It exposes the candidate row, normalization frontier, critical-pair ledger, endpoint rows, open socket, L10 exit certificate, and the structural $H,C,P,L$ rows as the finite handoff boundary.

Apply the candidate discharge schedule exhaustion. The consumer must read the bounded candidates together with the displayed discharge schedule before any conversion-facing export is formed.

Use the conditional budget and bounded checker gates. These gates keep the checker budget finite and tied to the same frontier and pair-ledger rows, so confluence-facing and decidability-facing consumers receive a bounded conversion surface rather than an unrestricted theorem.

Finally repack the local NameCert rows. The $\hsame$, $\Cont$, $\Pkg$, and naming coordinates preserve the accepted packet, while the socket and L10 rows prevent the route from becoming global strong normalization, unrestricted confluence, unbounded conversion decidability, quotient equality, selected limits, host equality, or Real completeness.
\end{proof}
\leanchecked{BEDC.Derived.MetaCICCriticalPathUp.MetaCICCriticalPathBoundedConversionConsumerRoute}

\begin{theorem}[MetaCIC critical-path L10 object-status readback]
\label{thm:metacic-critical-path-l10-object-status-readback}
Any candidate-mediated strong-normalization discharge used as a roadmap-critical $\MetaCICCriticalPathUp$ packet may cite L10 object readiness only through the displayed status rows for $\DyadicRatCoreUp$, $\StreamNameUp$, $\RegSeqRatUp$, $\RealUp$, and $\RealCompletionExactBoundaryUp$. The readback factors through the four-object exit certificate, the Real L10 objectwise audit pullback, and the exact-boundary packet rows; it must not turn those rows into unconditional $\RealUp$ completeness, quotient real equality, selected host limits, closed subject reduction, host normalization, global confluence, or beta-conversion decidability.
\end{theorem}
\begin{proof}
Start with the candidate-mediated SN discharge theorem, \autoref{thm:metacic-critical-path-candidate-mediated-sn-discharge}. It keeps the discharge socket $D$ open and requires the candidate row, endpoint rows, handoff row, structural rows, and four L10 faces before any confluence-facing or decidability-facing budget is read.

Next read the L10 object status route. The four-object exit certificate of \autoref{thm:metacic-critical-path-l10-four-object-exit-certificate} displays the $\DyadicRatCoreUp$, $\StreamNameUp$, $\RegSeqRatUp$, and $\RealUp$ faces, while \autoref{thm:metacic-critical-path-l10-status-bridge-obligation} and \autoref{thm:real-l10-current-phase-exit-audit-pullback} force those faces to contribute their own local closurestatus, formalstatus, leantarget, and bridgestatus rows.

Finally attach the completion boundary only as a finite exact-boundary read. The $\RealCompletionExactBoundaryUp$ packet supplies the displayed exact-boundary rows already governed by its local NameCert surface, so the SN-facing obligation may cite that boundary as status readback but cannot export Real completeness, host normalization, global confluence, quotient equality, selected limits, or a filled discharge socket.
\end{proof}
\leanchecked{BEDC.Derived.MetaCICCriticalPathUp.MetaCICCriticalPathL10ObjectStatusReadback}
